\section{Overview}
\label{sec:overview}

Artificial intelligence systems are undergoing a fundamental architectural shift. For years, AI models served as passive, on-demand tools: a developer provided a list of available functions alongside a user request, and the model was expected to select and invoke the correct function from that flat list. This paradigm, commonly referred to as \textit{tool use}, has proven fragile in practice. When dozens or hundreds of capabilities are loaded simultaneously into an agent's context window, the model's reasoning becomes diluted, hallucinated tool calls increase, and the cost of operation rises sharply with every additional capability exposed \cite{schick2023toolformer, patil2023gorilla}.

A newer architectural approach---now implemented in production autonomous agent frameworks such as OpenClaw \cite{openclaw_2025}, LangGraph \cite{langgraph_2024}, and emerging sovereign agent systems---replaces flat tool lists with a hierarchical capability structure known as \textit{agentic skills}. Rather than loading all functionality at once, agents interact with their capabilities through a three-tiered process called \textit{Progressive Disclosure}: first seeing only skill names and brief descriptions (the metadata tier), then loading the full procedural instructions of a selected skill when needed (the instructional tier), and finally executing the underlying scripts or external integrations only at the moment of action (the execution tier). This approach is designed to preserve the agent's reasoning capacity, reduce unnecessary token expenditure, and enforce explicit behavioral constraints during execution \cite{wang2024executable, openclaw_2025}.

Despite growing adoption of this architectural pattern in open-source and enterprise autonomous systems, there is currently no empirical study that rigorously measures whether progressive disclosure actually improves agent performance compared to traditional tool use. Questions remain open: Does the three-tier loading model reduce hallucinated tool calls or merely redistribute them? Does it make execution trajectories more compliant with defined policies? At what practical scale do the token efficiency gains become meaningful? This thesis answers these questions through a controlled empirical comparison between the two paradigms, using real autonomous agent frameworks, measurable trajectory metrics, and a standardized task benchmark.

\newpage
